% =====================================================================================================
% DRAFT for paper1.tex (not yet in the manuscript): new Sec. III subsection "Fourth and sixth order"
% and an addendum to Appendix A. REVTeX 4.2 / PRB style; uses the paper's macros and labels
% (eq:e0, eq:h, eq:e2, eq:F, eq:wedge, tab:h, app:pt). New labels: sec:pt46, eq:series, eq:eps4,
% eq:expand, eq:W, eq:w23, eq:width23, eq:width6, tab:eps, eq:ab.
%
% Every number carries a "% src:" comment. Sources (vu_work/10_results/bilayer_annni_paper_results/
% reanalysis_v3/perturbation/):
%   A4 = pt_order4_analysis.json   H4 = pt_order4_hull.json   V4 = pt_order4_verify.json
%   A6 = pt_order6_analysis.json   V6 = pt_order6_verify.json
%   T  = audit scratch pt_table.json, E6 = audit scratch pt6_edges.json. Both are produced by
%        pt_table.py and pt6_edges.py, which call only the repo functions per_column, coeffs,
%        per_spin_series and derivs_at_half. Before submission, fold them into pt_order6_v3.py --analyze
%        (or copy them to audit_9b/) so that T and E6 also live in the repo outputs.
%
% Suggested edits elsewhere (not made here):
%   * Sec. III.B, last paragraph: replace the \pending{Decide whether to carry the expansion to order g^4 ...}
%     and "higher orders decide whether mixed structures such as <2^k3> appear there" by a pointer to
%     Sec. \ref{sec:pt46}.
%   * Sec. III.B: the <3> wedge boundaries could quote the g^4 terms of Eq. (eq:w23) and of kappa_{FM|3}.
%   * Intro/abstract/conclusion: add the <23> result only after the S5a QMC analysis (plan step 5).
%   * refs: the \pending on Fisher-Selke/Fisher-Szpilka below needs a verified statement and bib entries
%     (Fisher & Szpilka, PRB 36, 644 and 5343 (1987) are verified per the project notes, but not yet in refs.bib).
% =====================================================================================================

\subsection{Fourth and sixth order}
\label{sec:pt46}

At $\kappa_{3|2}$ every sequence of two- and three-site domains has the same energy at order $g^2$, so
higher orders decide which structures are stable there. We carry the expansion to order $g^6$. Moving a
wall takes a whole row of flips, so structures with different walls mix only at order $g^{L_x}$
($g^{2L_x}$ for $r>0$). For $L_x\to\infty$ each straight-wall structure is therefore a nondegenerate
reference state, and Rayleigh--Schr\"odinger theory applies to each one separately. We write its energy
per spin as
\begin{equation}
e=e_0+\varepsilon_2g^2+\varepsilon_4g^4+\varepsilon_6g^6+O(g^8),
\label{eq:series}
\end{equation}
where $\varepsilon_2g^2$ is the term $e_2$ of Eq.~\eqref{eq:e2}. Odd orders vanish because every flipped
spin must be flipped back. Write the Ising part of Eq.~\eqref{eq:H} as
$-\sum J_{ij}\hat\sigma^z_i\hat\sigma^z_j$, so that $J_{ij}=-\kappa$ for the pairs at distance 2 along $y$.
Let $D_i=2h_i$ be the cost of flipping spin $i$, and let $B_{ij}=J_{ij}s_is_j$, which is minus the energy of
the bond between spins $i$ and $j$. Then
\begin{equation}
N\varepsilon_4=\sum_i\frac{1}{D_i^3}
-\sum_{\langle ij\rangle}\frac{4B_{ij}\,(D_i+D_j)}{(D_iD_j)^2\,(D_i+D_j-4B_{ij})} .
\label{eq:eps4}
\end{equation}
The first term is the $g^4$ term of an isolated spin, whose energy is $-(h^2+g^2)^{1/2}$. The second sum
runs over coupled pairs only. For an uncoupled pair, the fourth-order processes that flip both spins cancel
exactly against the normalization term of the expansion. This cancellation generalizes. Each spin
that takes part in a process is flipped and restored, so the $g^{2n}$ coefficient receives contributions
only from connected clusters of at most $n$ spins. We obtain $\varepsilon_6$ from this linked-cluster form,
solving every connected cluster of up to three spins exactly in rational arithmetic.
%
Equation~\eqref{eq:eps4} agrees with exact diagonalization of small clusters to $3\times10^{-6}$ relative.
% src: V4 closed_vs_ED, max |E4_ED/E4 - 1| = 2.7e-6
The sixth-order coefficient agrees with exact diagonalization to $10^{-4}$ relative,
% src: V6 vs_ED, max |rel| = 1.1e-4 (limited by the ED fit)
and with an independent evaluation that does not use the cluster decomposition to $10^{-15}$.
% src: audit scratch cmp_e6.py (13 structures), copied to perturbation/audit_9b/
Table~\ref{tab:eps} lists the coefficients at $\kappa=1/2$.

\begin{table}
\caption{Coefficients of Eq.~\eqref{eq:series} per spin at $\kappa=1/2$, in units of $J_0$, for the
structures that compete near $\kappa_{3|2}$; $\rho$ is the wall density. $W_4=k\varepsilon_2'+\varepsilon_4$
with $k=F(3)/4$ [Eq.~\eqref{eq:W}] is the fourth-order energy on the $\langle3\rangle|\langle2\rangle$ line
(not needed for FM). $\varepsilon_2$ and $\varepsilon_4$ are exact; $W_4$ and $\varepsilon_6$ are rounded.}
\label{tab:eps}
\begin{ruledtabular}
\begin{tabular}{llcccc}
$r$ & & $\rho$ & $\varepsilon_4$ & $10^3\,W_4$ & $10^5\,\varepsilon_6$ \\
\hline
0 & FM                    & 0   & $-5/1512$            & ---      & $-24.15$ \\
  & $\langle3\rangle$     & 1/3 & $-1/125$             & $-7.333$ & $-7.454$ \\
  & $\langle23\rangle$    & 2/5 & $-54673/7560000$     & $-4.610$ & $-120.5$ \\
  & $\langle2\rangle$     & 1/2 & $-13/4320$           & $2.546$  & $-22.06$ \\
1 & FM                    & 0   & $-1/1152$            & ---      & $-2.062$ \\
  & $\langle3\rangle$     & 1/3 & $-437/362880$        & $-1.011$ & $4.409$  \\
  & $\langle23\rangle$    & 2/5 & $-4477/3628800$      & $-0.5972$ & $-2.175$ \\
  & $\langle2\rangle$     & 1/2 & $-11/13440$          & $0.4836$ & $-2.055$ \\
\end{tabular}
\end{ruledtabular}
% src: T (all entries). FM/<2>/<3> eps4 also A4 exact_at_kappa_half; W4 of <3>, <23>, <2> also A4 hull_3_2.hull
% (exact: -11/1500, -34849/7560000, 11/4320 at r=0; -367/362880, -2167/3628800, 13/26880 at r=1).
% eps2 is not tabulated: it follows from Table I and Eq. (e2), e.g. -1/5, -14/75, -1/6 for <3>, <23>, <2> at r=0.
\end{table}

\paragraph*{Expansion about the multiphase point.}
Write $\kappa=\tfrac12+kg^2+k_1g^4+k_2g^6$. The classical energy~\eqref{eq:e0} is linear in $\kappa$ with
slope $1-4\rho$, so Eq.~\eqref{eq:series} becomes
\begin{align}
e={}&e_0\big|_{\kappa=1/2}+g^2\big[k(1-4\rho)+\varepsilon_2\big]\nonumber\\
&+g^4\big[k_1(1-4\rho)+W_4\big]\nonumber\\
&+g^6\big[k_2(1-4\rho)+W_6\big]+O(g^8),
\label{eq:expand}
\end{align}
with the $\kappa$ derivatives (primes) taken at $\kappa=1/2$:
\begin{equation}
W_4=k\varepsilon_2'+\varepsilon_4,\qquad
W_6=k_1\varepsilon_2'+\tfrac12k^2\varepsilon_2''+k\varepsilon_4'+\varepsilon_6 .
\label{eq:W}
\end{equation}
On the second-order boundary, $k=F(3)/4$, the $g^2$ bracket is the same for every sequence of two- and
three-site domains (Appendix~\ref{app:pt}). Each domain of $n\ge4$ sites adds $\tfrac12g^2[(n-2)F(3)-F(n)]>0$
to the numerator of Eq.~\eqref{eq:de}.
% src check: gives the minimum penalties 7/1320 (r=0) and 1/440 (r=1) of H4 min_g2_penalty_long_domains
At order $g^4$ only the sequences of 2s and 3s therefore compete, with energy
$k_1(1-4\rho)+W_4$ per spin. Since the first term is linear in $\rho$, the stable structures are the vertices
of the lower convex hull of $W_4$ plotted against $\rho$.

\paragraph*{Fourth order: the $\langle23\rangle$ phase.}
At $\kappa=1/2$ the field~\eqref{eq:h} of a spin depends only on its position in its own domain
(Table~\ref{tab:h}). The pairs in Eq.~\eqref{eq:eps4} lie either within one domain or on the two sides of a
single wall. $W_4$ is therefore a sum of terms for single domains and for pairs of adjacent domains
[Eq.~\eqref{eq:ab}]. A periodic sequence can be read as a closed path through the domain lengths, and its
energy per spin is the cost of the path divided by its length. Such a ratio is smallest on an elementary
cycle, so only $\langle2\rangle$, $\langle3\rangle$ and $\langle23\rangle$ ($\up\up\dn\dn\dn$) can be stable at this
order, among sequences of any length (Appendix~\ref{app:pt}). Table~\ref{tab:eps} places $\langle23\rangle$
below the chord from $\langle3\rangle$ to $\langle2\rangle$, by $619/504000$ per spin at $r=0$ and $667/3628800$
at $r=1$.
% src: H4 by_r.{0,1}.chord_depth_23 = -619/504000, -667/3628800
$\langle23\rangle$ is therefore stable in a window
\begin{equation}
\tfrac12+\tfrac14F(3)g^2+k_1^-g^4<\kappa<\tfrac12+\tfrac14F(3)g^2+k_1^+g^4,
\label{eq:w23}
\end{equation}
with $k_1^-=20591/2016000\approx0.0102$ and $k_1^+=6011/336000\approx0.0179$ at $r=0$, and
$k_1^-=167/107520\approx0.00155$ and $k_1^+=1961/725760\approx0.00270$ at $r=1$.
% src: H4 by_r.{0,1}.window_k1 (also A6 boundaries.*.k1)
The width is
\begin{align}
\Delta\kappa_{23}&=\tfrac{619}{80640}\,g^4\approx0.00768\,g^4 &&(r=0),\nonumber\\
&=\tfrac{667}{580608}\,g^4\approx0.00115\,g^4 &&(r=1).
\label{eq:width23}
\end{align}
% src: H4 by_r.{0,1}.window_width_g4; ratio H4 width_ratio_r0_r1 = 6.68
The bilayer window is 6.7 times narrower. This is a much stronger effect of the interlayer coupling than the
factor 2.4 for the $\langle3\rangle$ wedge. The window keeps shrinking with $r$: its width is $0.00274\,g^4$ at
$r=1/2$ and $0.00028\,g^4$ at $r=2$.
% src: H4 by_r.{1/2,2}.window_width_g4_float
The metastable branches $\langle3\rangle$ and $\langle2\rangle$ cross inside the window, at
$\kappa=\tfrac12+\tfrac14F(3)g^2+k_1^{32}g^4$ with $k_1^{32}=1067/72000$ ($r=0$) and $31/13824$ ($r=1$). This
is the line along which template free energies of $\langle3\rangle$ and $\langle2\rangle$ can be compared with
simulations.
% src: A4 boundaries.{r=0,r=1}.kappa_32
On the ferromagnetic side the boundary moves to
$\kappa_{\rm FM|3}=\tfrac12-\tfrac18F(3)g^2+k_1^{\rm F}g^4$, with $k_1^{\rm F}=-59/10080$ ($r=0$) and
$-1573/1935360$ ($r=1$).
% src: A4 boundaries.{r=0,r=1}.kappa_FM3
No new structure appears there. At $\kappa_{\rm FM|3}$ every structure other than the ferromagnet and
$\langle3\rangle$ contains two-site domains or domains of four or more sites. Each of these costs energy at
order $g^2$, which higher orders cannot overturn at small $g$.

\paragraph*{Degenerate edges.}
The argument also shows exactly which structures tie on the window edges. In the $(\rho,W_4)$ plane, a
sequence with no two adjacent three-site domains ($\langle223\rangle$, $\langle2223\rangle$,
$\langle22323\rangle,\dots$) lies exactly on the segment joining $\langle2\rangle$ and $\langle23\rangle$. A
sequence with no two adjacent two-site domains ($\langle233\rangle$, $\langle2333\rangle,\dots$) lies exactly on
the segment joining $\langle23\rangle$ and $\langle3\rangle$ (Appendix~\ref{app:pt}).
% src: H4 by_r.*.named_heights (all exactly 0)
Each edge of the window is therefore itself a multiphase line at order $g^4$, and the next order decides it.

\paragraph*{Sixth order.}
A connected cluster of three spins spans at most one intermediate domain. $W_6$ is therefore a sum of terms
for three consecutive domains, which we have confirmed in exact arithmetic.
% src: audit scratch e6_triples.py / e6_triples8.out (copied to audit_9b/): zero residual, 71 sequences, r=0,1
Applied to pairs of consecutive domains, the elementary-cycle argument allows at most one new structure in
each gap at this order: $\langle223\rangle$ between $\langle23\rangle$ and $\langle2\rangle$, or
$\langle233\rangle$ between $\langle3\rangle$ and $\langle23\rangle$. $\langle223\rangle$ opens, in a window of
width $4.1\times10^{-4}g^6$ at $r=0$ and $2.3\times10^{-5}g^6$ at $r=1$.
% src: A6 boundaries.*.23|2.hull_at_g6 = [23, 223, 2]; widths E6 width_223_g6 = 4.146e-4, 2.301e-5
% (depths of <223> below the <23>-<2> chord: A6 -3.384e-5, -1.878e-6)
$\langle233\rangle$ does not open. It lies above the segment from $\langle3\rangle$ to $\langle23\rangle$ by
$2.0\times10^{-6}g^6$ per spin at $r=0$ ($1.2\times10^{-7}g^6$ at $r=1$), a margin small enough that the next
order could reverse it.
% src: A6 boundaries.*.3|23.hull_at_g6 = [3, 23]; heights E6 height_233_vs_chord_3_23 = +1.97e-6, +1.24e-7
Through order $g^6$ the transverse field thus orders the structures as $\langle3\rangle$, $\langle23\rangle$,
$\langle223\rangle$, $\langle2\rangle$ with increasing $\kappa$. We do not follow the sequence further.
The same order of appearance occurs in the low-temperature expansion of the classical three-dimensional
model~\cite{fisher1980infinitely,selke1988annni}. There $\langle3\rangle$ appears at first order with a
first-order boundary to the ferromagnet. $\langle23\rangle$ springs from the multiphase point at second
order, with a first-order $\langle3\rangle|\langle23\rangle$ boundary. The phases $\langle2^k3\rangle$
follow on the $\langle2\rangle$ side at successively higher orders. In the language of domain-wall
interactions~\cite{fisher1987domainI,fisher1987domainII}, $g^2$ takes over the role of the low-temperature
expansion parameter. The selecting excitations are virtual rather than thermal spin flips next to the walls,
and the coefficients are different. The quantum result does not follow from the classical one: at $T=0$ the
model maps onto a strongly anisotropic three-dimensional classical model~\cite{suzuki1976relationship}, which is
not in the regime of the classical low-temperature expansion.
% src: read in Selke, Phys. Rep. 170, 213 (1988), Sec. 4.2 (2026-09-29). fisher1987domainI/II = Fisher & Szpilka,
% PRB 36, 644 and 5343 (1987), verified via Crossref; bib entries still to be added to refs.bib.

The sixth order also shifts the edges of the $\langle23\rangle$ window. With
$k_2^-=-3.06\times10^{-3}$ and $k_2^+=+5.26\times10^{-3}$ at $r=0$, and $-1.43\times10^{-4}$ and
$+2.42\times10^{-4}$ at $r=1$ (the upper edge is now the $\langle23\rangle|\langle223\rangle$ boundary),
% src: E6 k2_3|23, k2_23|223
\begin{align}
\Delta\kappa_{23}&=0.00768\,g^4\,(1+1.08\,g^2) &&(r=0),\nonumber\\
&=0.00115\,g^4\,(1+0.34\,g^2) &&(r=1).
\label{eq:width6}
\end{align}
% src: E6 width_23_rel_g6_over_g4 = 1.085, 0.335
The correction is 27\% of the leading term at $g=0.5$ for $r=0$, and it equals the leading term near $g=1$.
Equation~\eqref{eq:width23} is therefore quantitative only at small $g$. Where the window becomes wide enough
to resolve numerically, the sixth-order terms are needed, and at $r=0$ the series converges slowly there.

% =====================================================================================================
% Addendum to Appendix A ("Minimum of the domain-sequence energy"), to follow its last paragraph.
% =====================================================================================================

\paragraph*{Higher orders.}
At order $g^4$ on the line $k=F(3)/4$ the same ratio structure appears, with the domain costs replaced by
costs of adjacent domain pairs. Let $P(D,D',B)=-4B(D+D')/[(DD')^2(D+D'-4B)]$ be the pair term of
Eq.~\eqref{eq:eps4}. Let $D_1,\dots,D_n$ be the flip costs along a domain (Table~\ref{tab:h}), and let
$\gamma_i=s_i(s_{i+2}+s_{i-2})$, which takes the values $-2$, $0$ and $2$. Let $X_n$ and $Y_n$ be the flip costs of
the first and second site from a wall of a domain of length $n$. Then, for a period of $L$ sites,
\begin{align}
W_4L&=\sum_k\big[a(n_k)+b(n_k,n_{k+1})\big],\nonumber\\
a(n)&=\sum_{i=1}^{n}\Big[-\frac{2k\gamma_i}{D_i^2}+\frac1{D_i^3}\Big]\nonumber\\
&\quad+\sum_{i=1}^{n}\Big[P(D_i,D_i,1)+\tfrac12P(D_i,D_i,r)\Big]\nonumber\\
&\quad+\sum_{i=1}^{n-1}P(D_i,D_{i+1},1)+\sum_{i=1}^{n-2}P(D_i,D_{i+2},-\tfrac12),\nonumber\\
b(n,n')&=P(X_n,X_{n'},-1)+P(Y_n,X_{n'},\tfrac12)\nonumber\\
&\quad+P(X_n,Y_{n'},\tfrac12).
\label{eq:ab}
\end{align}
% src: audit scratch domain_form.py (copied to audit_9b/): reproduces W4 exactly for 1167 sequences, r = 0, 1
With $C(n,n')=k_1(n-4)+a(n)+b(n,n')$, the fourth-order energy per spin is $\sum_kC(n_k,n_{k+1})/\sum_kn_k$.
A periodic sequence is a closed walk on the directed graph with vertices $\{2,3\}$ and all four edges, where
edge $n\to n'$ has cost $C(n,n')$ and length $n$. Every closed walk splits into elementary cycles. A ratio of
sums with positive denominators is at least the smallest ratio of its parts, with equality only if all parts
share that ratio. The elementary cycles are $2\to2$, $3\to3$ and $2\to3\to2$, that is $\langle2\rangle$,
$\langle3\rangle$ and $\langle23\rangle$, so one of them minimizes the energy for every $k_1$.

Now count the adjacent pairs $N_{22}$, $N_{33}$ and $N_{23}=N_{32}$ of a cyclic sequence. The period is
$L=2N_{22}+3N_{33}+5N_{23}$, and the total cost is $N_{22}C(2,2)+N_{33}C(3,3)+N_{23}[C(2,3)+C(3,2)]$. The point
$(\rho,W_4)$ of the sequence is therefore the average of the points of $\langle2\rangle$, $\langle3\rangle$ and
$\langle23\rangle$, weighted by the site fractions $2N_{22}/L$, $3N_{33}/L$ and $5N_{23}/L$. Sequences with
$N_{33}=0$ lie on the segment from $\langle2\rangle$ to $\langle23\rangle$, and sequences with $N_{22}=0$ lie on
the segment from $\langle23\rangle$ to $\langle3\rangle$. All other sequences lie strictly above the hull,
because $\langle23\rangle$ lies strictly below the chord from $\langle3\rangle$ to $\langle2\rangle$.

At order $g^6$ the terms of Eq.~\eqref{eq:W} involve at most three consecutive domains. The same argument,
applied to the graph whose vertices are pairs of consecutive domains, leaves $\langle2\rangle$,
$\langle23\rangle$ and $\langle223\rangle$ as the only candidates on the $\langle23\rangle|\langle2\rangle$
line, and $\langle3\rangle$, $\langle23\rangle$ and $\langle233\rangle$ on the $\langle3\rangle|\langle23\rangle$
line.
